% Root proof; full independent written audit by descent_groups, with clarifications.
\subsection{C.3: Galois indices over a perfect coefficient field}

Let $k$ be a perfect field of characteristic $p>2$, fix an algebraic
closure $\bar k$, and put $\Gamma=\operatorname{Gal}(\bar k/k)$.
Use the standard split Gram matrix with entries in $\mathbb F_p$.
Thus the coordinate Hopf algebra and its degree filtration are defined
over $\mathbb F_p$. Let $\mathscr I_{\bar k}$ be the explicit canonical
index set of Theorem~\ref{C:Om-canonical-templates}.
For $\gamma\in\Gamma$ define $\gamma\cdot i$ by applying $\gamma$
to the coefficient matrices of $U_i$ and then canonicalizing them.
This is an action on the index set even though its chosen total order
need not be Galois invariant. Each orbit is finite: the finitely many
matrix coefficients belong to a finite Galois extension of $k$.
Choose the least index of each finite orbit. These minima form an
explicit set $\mathscr O_k$; the equivalence being removed is the
specified Galois action on already classified matrices.

We distinguish a classification by these indices and finite-dimensional
algebra data from a uniform decision algorithm over arbitrary computable
fields. No algorithm for central division algebras or hermitian forms
over an arbitrary field is assumed without saying so.

\begin{lemma}[Proved: reflection of representation isomorphisms]
\label{C:NoetherDeuring_rational}
Let $k\subseteq L$ be any extension of fields and let $M,N$ be finite
dimensional rational representations of an affine group scheme over $k$.
If $M\otimes_kL\simeq N\otimes_kL$, then $M\simeq N$.
For infinite $k$, an isomorphism can be found by a finite determinant
grid after solving linear intertwiner equations over $k$.
\end{lemma}
\begin{proof}
The coaction coefficients of $M\oplus N$ lie in a finite-dimensional
coalgebra. Its dual is a finite-dimensional $k$-algebra, and the
assertion is the Noether--Deuring theorem for its finite-dimensional
modules, in \cite[Theorem 9.2.8]{C-Webb2016}.
For the algebraic extensions used in the descent argument below,
there is also a direct proof. An isomorphism and its inverse use
only finitely many coefficients, so are defined over a finite
intermediate extension $L_0/k$. Forgetting $L_0$-scalars yields
$M^{\oplus[L_0:k]}\simeq N^{\oplus[L_0:k]}$ over $k$.
Krull--Schmidt uniqueness cancels this common positive multiplicity.
This argument does not require $k$ to be infinite or $L_0/k$ separable.
Here is the elementary constructive argument for the infinite-field
case used below. The vector space $V=\operatorname{Hom}_G(M,N)$ is
the kernel of finitely many linear equations in matrix entries.
Tensoring with $L$ therefore gives
$L\otimes_k V=\operatorname{Hom}_{G_L}(M_L,N_L)$.
For a $k$-basis $f_1,\ldots,f_s$ of $V$, the polynomial
$\det(\sum_j z_jf_j)$ has a nonzero value over $L$, hence is not
the zero polynomial over $k$. If $n=\dim_k M=\dim_k N$, a grid
of $n+1$ distinct elements of $k$ in every variable contains a
nonzero value, by the degree bound in each variable. The corresponding
linear combination is the desired equivariant isomorphism.
\end{proof}

\begin{theorem}[Proved: complete indecomposable Galois indices]
\label{C:perfect_Galois_indices}
For every perfect $k$ as above, $\mathscr O_k$ is in bijection with
the indecomposable finite-dimensional rational $O_{m,k}$-modules.
Write $M_\Omega$ for the module attached to the orbit $\Omega$.
There is a positive integer $s_\Omega$ such that
\[
 (M_\Omega)_{\bar k}
       \simeq\bigoplus_{i\in\Omega}U_i^{\oplus s_\Omega}.
                                                        \tag{C.3.G1}
\]
The integer is the degree of the residue division algebra
$D_\Omega=\operatorname{End}(M_\Omega)/\operatorname{rad}$ over
its center $K_\Omega$:
\[
 [K_\Omega:k]=|\Omega|,\qquad
 s_\Omega^2=\dim_{K_\Omega}D_\Omega.
                                                        \tag{C.3.G2}
\]
All this data can be obtained from one explicitly constructed
finite-dimensional endomorphism algebra over $k$: choose a finite
Galois extension $L/k$ defining one $U_i$, form the restriction of
scalars $W=\operatorname{Res}_{L/k}U_i$, and compute
$E=\operatorname{End}_{O_{m,k}}(W)$ by linear equations.
Then $E/\operatorname{rad}E\simeq M_r(D_\Omega)$ for some $r$,
and the image on $W$ of any lifted primitive idempotent of this
simple algebra is $M_\Omega$. Different choices give isomorphic
modules and the same index and division-algebra isomorphism class.
The least coefficient degree is constant on $\Omega$ and is also
the least coefficient degree of $M_\Omega$.
\end{theorem}
\begin{proof}
We first analyze any indecomposable $M$ over $k$. Its endomorphism
algebra $E_M$ is finite dimensional and local, by the Fitting lemma.
Let $J_M$ be its nilpotent radical, $D_M=E_M/J_M$, and let $K_M$
be the center of the division algebra $D_M$. Since $k$ is perfect,
$K_M/k$ is finite separable. Central-simple splitting
\cite[Lemma 11.8.6, Tag 0753]{C-StacksPerfectDescent} gives
\[
 D_M\otimes_k\bar k
   \simeq\prod_{\tau:K_M\hookrightarrow\bar k}M_s(\bar k),
       \qquad s^2=\dim_{K_M}D_M.
                                                        \tag{C.3.G3}
\]
The ideal $J_M\otimes_k\bar k$ is nilpotent, and the quotient
in (C.3.G3) is semisimple. Therefore it is exactly the radical
of $E_M\otimes_k\bar k$. Linear intertwiner equations identify
that algebra with $\operatorname{End}(M_{\bar k})$.

Decompose $M_{\bar k}$ into indecomposables. Its semisimple
endomorphism quotient is the product of matrix algebras whose
sizes are their multiplicities; every absolute indecomposable
has residue endomorphism field $\bar k$.
Comparing with (C.3.G3) proves that there are $[K_M:k]$ distinct
absolute indecomposable classes, each with multiplicity $s$.
The Galois action on the factors in (C.3.G3) is transitive,
because it is the action on the embeddings of the field $K_M$.
The factors correspond to the isotypic classes of $M_{\bar k}$,
so these classes form exactly one orbit in $\mathscr I_{\bar k}$.
This proves (C.3.G1)--(C.3.G2) for every existing $M$.

Conversely fix an orbit $\Omega$ and a representative $U_i$.
Its finitely many coefficient matrices are defined over a finite
Galois extension $L/k$. Call that representation $U$ over $L$;
it is absolutely indecomposable, since its extension to $\bar k$
is $U_i$. Restriction of scalars is an actual rational representation
$W$ over $k$: in an $L/k$ basis replace each coefficient by its
multiplication matrix. Separability gives the decomposition
\[
 W_{\bar k}\simeq
       \bigoplus_{\sigma:L\hookrightarrow\bar k}{}^\sigma U.
                                                        \tag{C.3.G4}
\]
The support in (C.3.G4) is exactly $\Omega$, with constant
multiplicity $[L:k]/|\Omega|$.

Let $M$ and $N$ be two indecomposable summands of $W$.
By the preceding paragraph each has geometric support one
Galois orbit; since their supports are contained in $\Omega$,
they both have support all of $\Omega$.
Let their geometric multiplicities be $s,t$, respectively.
Then $(M^{\oplus t})_{\bar k}\simeq(N^{\oplus s})_{\bar k}$.
Lemma~\ref{C:NoetherDeuring_rational} reflects this isomorphism
to $k$. Krull--Schmidt uniqueness implies $M\simeq N$ and
$s=t$. Thus $W$ is isotypic, $W\simeq M_\Omega^{\oplus r}$.
Every other indecomposable with support $\Omega$ is isomorphic
to this one by the same comparison. This proves both existence
and uniqueness, without assuming that $U_i$ itself descends
to its field of moduli.

The finite-dimensional algebra $E$ is computed by the displayed
linear intertwiner equations, including the group-scheme coaction
identities, rather than just the action on rational points.
Its quotient is $M_r(D_\Omega)$ by the isotypic decomposition.
Primitive idempotents lift across its nilpotent radical.
Their images in $W$ are indecomposable and are all isomorphic
by what we have just proved. This supplies the promised
construction and proves its independence of $L$ and of all
idempotent choices.

Finally, the coefficient filtration is defined over $\mathbb F_p$,
so its least degree is unchanged by Galois conjugacy. Restriction
of scalars and passage to a direct summand use constant matrices
and do not increase this degree. Conversely, faithful scalar
extension detects membership of every coefficient in a filtration
piece. Formula (C.3.G1) therefore shows that the least degree of
$M_\Omega$ equals the common least degree of its geometric summands.
\end{proof}

\begin{corollary}[Proved: the split Green ring as an explicit Galois subring]
\label{C:perfect_Green_image}
Let $A_k(O_m)$ denote the split Green ring of finite-dimensional
rational representations. Extension to $\bar k$ is injective,
and its image is the free abelian subgroup with basis
\[
                s_\Omega\sum_{i\in\Omega}[U_i]
                        \quad(\Omega\in\mathscr O_k).
                                                        \tag{C.3.G5}
\]
This subgroup is closed under tensor products and all exterior
operations. Its products and exterior operations are obtained
from the algorithms over $\bar k$: the coefficient of $M_\Psi$
is the coefficient of any $U_j$, $j\in\Psi$, divided by
$s_\Psi$. In particular the division is always integral.
The dual index is the orbit obtained by dualizing the indices
in $\Omega$; its Schur index equals $s_\Omega$.
\end{corollary}
\begin{proof}
The split Green group has its Krull--Schmidt basis
$[M_\Omega]$. Their images (C.3.G5) have pairwise disjoint
supports, proving injectivity and the exact description of the
image. Scalar extension commutes with actual tensor products
and exterior powers, hence the subgroup is closed under them.
Expanding the image of an operation in this basis gives exactly
the coefficient rule and its integrality. The finite support
of every tensor or exterior operation makes each coefficient
calculation finite after the indicated arithmetic data are given.
The same reasoning applied to duals, which commute with scalar
extension, proves the last assertion. Virtual operations follow
coefficientwise from $\lambda_t(-x)=\lambda_t(x)^{-1}$.
\end{proof}

\paragraph{What this does and does not compute.}
The theorem replaces the unknown modular indecomposable classes over
a perfect field by the specified finite Galois orbits of the already
explicit algebraically closed templates. It retains genuine arithmetic
data: the division part and index of a specified finite-dimensional
algebra. For the split orthogonal group, Proposition~\ref{C:field_WG_structure}
then attaches the hermitian group of $D_\Omega$ with the adjoint
involution to each self-dual orbit, and a hyperbolic integer to each
other dual pair. An actual representative is obtained by the primitive
idempotent construction. Its dual and invariant bilinear forms are
linear systems; a nonsingular form of one of the two signs exists
by the checked hermitian-category theorem.

Over a general perfect field, forgetting forms need not be injective.
Thus the coefficient rule (C.3.G5) alone is not a rule for all
orthogonal-form coefficients: the residual hermitian multiplicity
forms must also be computed. We do not call the arithmetic
classification of these forms, or effective division-algebra splitting
over an arbitrary field, completed by this corollary. The perfectness
hypothesis is also used essentially in (C.3.G3); purely inseparable
centers over an imperfect field are not silently replaced by
separable ones.

The failure of an unrestricted orbit-existence argument over an
imperfect field can already be seen for a finite-dimensional
coefficient algebra. Take $k=\mathbb F_p(a)$ and
$B=k[t]/(t^p-a)$. Over $\bar k$ the one-dimensional module on
which $t$ acts as $\alpha=a^{1/p}$ has a singleton Galois orbit.
No nonzero $B$-module over $k$ becomes a direct sum of copies of
that module: such a scalar extension would make the matrix of $t$
equal to $\alpha I$, which is impossible for a matrix over $k$.
Instead, restriction of scalars from $k(\alpha)$ becomes a
length-$p$ Jordan module. This is a boundary example for the
general algebraic descent argument; it is not asserted here to
be a new orthogonal-group counterexample.
